\documentclass[11pt,a4paper]{article}
% Preámbulo compartido para los modelos del segundo parcial de ACN.
% Replica el formato del parcial real 2P-2S-2025.
\usepackage[utf8]{inputenc}
\usepackage[T1]{fontenc}
\usepackage[spanish]{babel}
\usepackage{amsmath,amssymb}
\usepackage[margin=2.6cm]{geometry}
\usepackage{enumitem}
\usepackage{array}

\pagestyle{plain}
\setlength{\parindent}{0pt}
\setlength{\parskip}{5pt}

% Listas de incisos: a), b), c), ...
\setlist[enumerate,1]{label=\alph*),leftmargin=2.2em,itemsep=7pt,topsep=7pt}
% Listas numeradas 1., 2., 3., ...
\newlist{puntos}{enumerate}{1}
\setlist[puntos]{label=\arabic*.,leftmargin=2.2em,itemsep=7pt,topsep=7pt}
% Listas de condiciones del enunciado
\setlist[itemize,1]{label=$\bullet$,leftmargin=2em,itemsep=5pt,topsep=6pt}

% Encabezado del parcial
\newcommand{\encabezado}[1]{%
  \begin{center}
    \textbf{\large Aplicaciones Computacionales en Negocios}\\[3pt]
    \textbf{\large Segundo semestre de 2025 --- Segundo parcial}\\[3pt]
    \textbf{\normalsize Modelo de práctica #1}
  \end{center}
  \vspace{4pt}
  \hrule height 0.8pt
  \vspace{14pt}
  Nombre:\par\vspace{12pt}
  Legajo:\par\vspace{12pt}
  DNI:\par\vspace{12pt}
  ¿Te recibís con este parcial?:
  \vspace{10pt}
  \hrule height 0.8pt
  \vspace{16pt}
}

% Título de ejercicio
\newcommand{\ejercicio}[1]{\par\medskip\noindent\textbf{Ejercicio #1.}\ }

% Atajos de notación usados por la cátedra
\newcommand{\R}{\mathbb{R}}
\newcommand{\Z}{\mathbb{Z}}
\newcommand{\dem}{\mathrm{dem}}
\newcommand{\capa}{\mathrm{cap}}
\newcommand{\costo}{\mathrm{costo}}


\begin{document}
\encabezado{2}

%======================================================================
\ejercicio{1}
Consideremos el siguiente modelo de \emph{facility location}, en el que hay un
conjunto $C$ de clientes y un conjunto $L$ de ubicaciones candidatas:

\begin{align}
\min \quad & \sum_{i \in C} \sum_{j \in L} \alpha\, d_{ij}\, x_{ij}
           + \sum_{j \in L} \beta\, y_j \\
\sum_{j \in L} x_{ij} &= 1 && \forall\, i \in C \\
x_{ij} &\le y_j && \forall\, i \in C,\ \forall\, j \in L \\
x_{ij} &\in \{0,1\} && \forall\, i \in C,\ \forall\, j \in L \notag \\
y_j &\in \{0,1\} && \forall\, j \in L \notag
\end{align}

En este modelo, la variable binaria $x_{ij}$ representa si el cliente $i$ es
atendido desde la ubicación $j$, y la variable binaria $y_j$ especifica si se
abre un \emph{facility} en la ubicación $j$. El parámetro $d_{ij} \in \R_+$ es la
distancia entre el cliente $i$ y la ubicación $j$, la constante $\alpha > 0$ es el
costo por kilómetro de viaje a lo largo del período de operación, y la constante
$\beta > 0$ es el costo de instalación de cada \emph{facility}.

\begin{enumerate}
  \item ¿Cuál es la solución óptima de este modelo si $\beta = 0$? ¿Y si además
  eliminamos la restricción (3)? Justificar en ambos casos.

  \item Una formulación alternativa reemplaza la restricción (3) por
  \[
  \sum_{i \in C} x_{ij} \;\le\; |C|\; y_j \qquad \forall\, j \in L .
  \]
  Explicar por qué ambas formulaciones tienen exactamente las mismas soluciones
  enteras factibles. ¿Son igualmente buenas a la hora de resolver el modelo con
  un \emph{solver}? Justificar exhibiendo una solución de la relajación lineal
  que sea factible para una de las formulaciones y no para la otra.

  \item Supongamos que cada cliente $i \in C$ tiene una demanda
  $\dem_i \in \R_+$ asociada y que cada ubicación $j \in L$ tiene una capacidad
  máxima $\capa_j \in \R_+$. ¿Cómo cambia el modelo si necesitamos agregar esta
  restricción? Supongamos además que no queremos abrir un \emph{facility} si la
  suma de las demandas de los clientes asignados a ese \emph{facility} es menor
  a un umbral $\delta$. ¿Cómo se agrega esta restricción al modelo?

  \item Los clientes están agrupados en cuadras: para cada $i \in C$ tenemos el
  dato $\mathrm{cuadra}_i \in \Z$. Agregar al modelo una restricción para que
  dos clientes de la misma cuadra sean siempre atendidos desde la misma
  ubicación. ¿Cuántas restricciones se agregan al modelo con esta condición?
  ¿Se puede lograr el mismo efecto sin agregar restricciones?
\end{enumerate}

\newpage
%======================================================================
\ejercicio{2}
Tenemos que diseñar el \emph{fixture} de una liga deportiva. Tenemos un conjunto
$E = \{1,\dots,n\}$ de equipos ($n$ es par) y un conjunto
$F = \{1,\dots,2n-2\}$ de fechas. Los equipos deben jugar todos contra todos,
ida y vuelta (\emph{double round-robin}). Para cada par de equipos
$i,j \in E$, $i \ne j$, y cada fecha $f \in F$, consideramos la variable binaria
$x_{ijf}$, que toma el valor 1 si el equipo $i$ juega de local contra el equipo
$j$ en la fecha $f$, y 0 en caso contrario.

\begin{puntos}
  \item Escribir las restricciones que garantizan que (i) cada equipo juega
  exactamente un partido en cada fecha, y (ii) cada partido se juega exactamente
  una vez a lo largo del campeonato. Observar que, con esta formulación, no hace
  falta una variable adicional para representar la localía: explicar por qué, y
  escribir la expresión que indica si el equipo $i$ juega de local en la fecha
  $f$.

  \item Agregar restricciones para que ningún equipo juegue más de dos partidos
  consecutivos de local ni más de dos partidos consecutivos de visitante.

  \item Sea $K \subseteq E \times E$ el conjunto de los pares de equipos que son
  rivales clásicos. Agregar restricciones para que los clásicos se jueguen entre
  las fechas 4 y 12, y para que el equipo 5, que no dispone de su estadio en las
  fechas 3 y 11, no juegue de local en esas fechas.

  \item Decimos que un equipo tiene un \emph{break} en la fecha $f > 1$ si juega
  con la misma localía en las fechas $f-1$ y $f$. Agregar al modelo las
  variables y restricciones necesarias para minimizar la cantidad total de
  \emph{breaks} del campeonato. Explicar por qué, con la función objetivo
  propuesta, alcanza con una sola desigualdad por variable de \emph{break} para
  que el modelo sea correcto, y por qué estas variables pueden declararse
  continuas en $[0,1]$ en lugar de binarias.

  \item Decimos que un equipo tiene un \emph{patrón perfecto} si alterna siempre
  local y visitante, es decir, si no tiene ningún \emph{break}. Explicar por qué
  a lo sumo dos equipos pueden tener un patrón perfecto en un mismo
  \emph{fixture}.
\end{puntos}

\newpage
%======================================================================
\ejercicio{3}
Somos responsables de la planificación de la producción de un laboratorio
medicinal para los próximos doce meses. Sea $P$ el conjunto de productos y sea
$T = \{1,\dots,12\}$ el conjunto de meses. Para cada producto $i \in P$ y cada
mes $t \in T$ tenemos la demanda estimada $d_{it} \in \R_+$, el costo unitario de
producción $c_{it} \in \R_+$ y el costo unitario de almacenamiento
$h_i \in \R_+$. El stock inicial de cada producto es $s_{i0} = 0$. La
planificación debe respetar las siguientes condiciones:

\begin{itemize}
  \item La planta puede fabricar como máximo 300 unidades por mes, sumando todos
  los productos.
  \item El depósito es compartido y no puede almacenar más de $S_{\max}$
  unidades en total al final de cada mes.
  \item La producción propia se realiza en lotes exactos de 10 unidades: en cada
  mes, de cada producto se fabrican 0, 10, 20, \dots\ unidades.
  \item Además de fabricar, se puede tercerizar producción a un proveedor
  externo, a un costo de $g$ pesos por unidad. El proveedor no puede entregar
  más de 200 unidades por mes sumando todos los productos y, si en un mes se le
  terceriza un producto, se le deben pedir al menos 20 unidades de ese producto.
  \item Cada producto tiene un stock de seguridad $m_i$ que se debe mantener al
  final de cada mes. Esta condición se puede incumplir, pero se paga una
  penalidad fija de $\Pi$ pesos por cada par (producto, mes) en el que se la
  incumpla.
\end{itemize}

\begin{puntos}
  \item Formular un modelo de programación lineal entera mixta para determinar
  cuánto fabricar y cuánto tercerizar de cada producto en cada mes, minimizando
  el costo total. Definir con precisión los conjuntos, los parámetros y las
  variables de decisión que utilice el modelo.

  \item Justificar por qué la penalidad del último punto requiere una variable
  binaria, mientras que una penalidad proporcional a las unidades faltantes no
  la requeriría.

  \item ¿Cómo se modifica el modelo si queremos incorporar una tasa de descuento
  $r$, de modo que los costos del mes $t$ se expresen en valor presente?

  \item El gerente observa que la solución óptima tiende a terminar el horizonte
  de planificación con stock nulo de todos los productos, lo cual deja a la
  planta en mala posición para el año siguiente. Explicar por qué ocurre esto y
  proponer una modificación del modelo que corrija este comportamiento.
\end{puntos}

\end{document}
